\section{Diseño de algoritmos de aprendizaje por refuerzo}

\subsection{PRIME: desde la perspectiva del Value Model}

Cuando el equipo de Cui Ganqu (崔干曲) intentó reproducir O1 a inicios de 2024 (antes de DeepSeek R1), el problema central que enfrentaron fue: \textbf{¿cómo superar la dificultad de Credit Assignment que provoca el Sparse Reward?}

Los métodos tradicionales (como GRPO) solo usan el Verifiable Reward final como señal de entrenamiento, lo que significa que el modelo solo puede obtener retroalimentación a partir del resultado final, sin poder saber con precisión qué pasos del proceso de razonamiento fueron correctos. La idea central de PRIME es introducir el \textbf{Implicit PRM} (modelo de recompensa de proceso implícito), que realiza el Credit Assignment de manera automática mediante la Log Probability.

\begin{importantbox}{La idea de diseño central de PRIME}
\begin{itemize}
    \item \textbf{Problema}: al inicio del entrenamiento, el Value Model de PPO predice muy mal en los modelos de lenguaje grandes, lo que provoca que la reward primero baje y luego se recupere (el fenómeno de «dip-then-recover»), introduciendo inestabilidad adicional en el entrenamiento
    \item \textbf{Solución}: modelar el Value mediante la Log Probability, en lugar de la forma tradicional de Linear Head. Antes y después del entrenamiento, el Reward Model sigue teniendo la forma de un modelo de lenguaje, lo que evita la práctica tosca de descartar la capa de embedding y conectarla a un linear head
    \item \textbf{Ventaja}: esta es una manera más natural y más robusta de modelar el Value, que en particular no presenta un sesgo grave de Value prediction al inicio del entrenamiento
\end{itemize}
\end{importantbox}

Visto desde la perspectiva actual, PRIME es en realidad una \textbf{forma alternativa de modelar el Value Model}. Aunque los métodos sin Value Model, como GRPO, ya logran entrenar muy bien, Cui Ganqu considera que, en escenarios futuros más complejos (como el Agentic RL con feedback interno), cómo convertir el feedback de múltiples fuentes en Process Reward o Value Estimation sigue siendo un problema que merece un estudio a fondo.

Además, PRIME tiene una conexión profunda con el método de Uncertainty Distillation: ambos minimizan la divergencia KL entre el Policy Model y el Teacher/Reward Model, lo cual es, en esencia, un proceso de destilar el conocimiento del Reward Model hacia el Policy Model.

\begin{knowledgebox}{El futuro de los métodos Value-based Actor-Critic}
Actualmente, métodos como GRPO predominan por su simplicidad y eficiencia, pero los métodos Value-based no están obsoletos. Cui Ganqu predice que, a medida que aumente la complejidad de las tareas (más fuentes de feedback, cadenas de interacción más largas), la comunidad volverá a prestar atención a los métodos Value-based Actor-Critic. Esta dirección todavía tiene un amplio espacio por explorar a fondo.
\end{knowledgebox}

\subsection{Reinforce++: el regreso a la sabiduría clásica}

El punto de partida de Hu Jian (胡建) al diseñar Reinforce++ tiene un matiz de «regreso a lo clásico». En la época tradicional del Multi-Agent RL, PPO mostraba un desempeño muy estable en escenarios de videojuegos a gran escala. Al llegar la era de los modelos grandes, la intuición de Hu Jian fue que muchos de los detalles de diseño de los métodos clásicos (como el Global Batch Norm) fueron cuidadosamente pensados por quienes los precedieron, y deberían heredarse en lugar de descartarse.

\begin{importantbox}{La filosofía de diseño de Reinforce++}
\begin{enumerate}
    \item Partir de PPO y \textbf{eliminar el Critic} (para evitar la dificultad de entrenar el Value Model en modelos grandes)
    \item \textbf{Conservar los trucos de ingeniería de PPO que ya han sido validados}: el Global Batch Norm, entre otros
    \item Se confirmó que los trucos clásicos siguen siendo efectivos en los nuevos escenarios: «lo que ha existido durante mucho tiempo, sin duda tiene una razón para existir»
\end{enumerate}
\end{importantbox}

Hu Jian planteó además un punto de vista importante: el algoritmo de RL del futuro ya no será un algoritmo aislado, sino \textbf{una combinación de un buen conjunto de Recipes}. Por ejemplo: Truncated Importance Sampling + On-Policy Learning + Batch Norm + una implementación de Infra adecuada, que en conjunto forman un esquema de entrenamiento estable.

\subsection{GSPO: el regreso de la optimización a nivel de Sequence}

El equipo de Zheng Chujie (郑楚杰) descubrió, durante el entrenamiento de la versión Qianwen 2507, que GRPO siempre presentaba problemas de estabilidad en las etapas finales. Sospecharon que la causa raíz estaba en \textbf{el objetivo de optimización mismo}: el objetivo de optimización de GRPO debería ser, por naturaleza, a nivel de Sequence, pero carece de una relación teórica directa con el objetivo de optimización a nivel de Token.

\begin{importantbox}{La innovación central de GSPO}
\begin{itemize}
    \item \textbf{Regreso a los primeros principios}: derivar directamente el objetivo de optimización a nivel de Sequence
    \item \textbf{Eliminación del KL Penalty}: en modelos de 235B y de mayor tamaño, eliminar el KL Penalty aumenta notablemente la velocidad de entrenamiento, y GSPO logra mantener un entrenamiento estable
    \item \textbf{Aplicación práctica}: GSPO se aplicó al lanzamiento de Qianwen 2507 y al entrenamiento de modelos posteriores (como Consonant X)
\end{itemize}
\end{importantbox}

\begin{warningbox}{Las limitaciones de GSPO}
La versión inicial de GSPO no consideraba el problema de la inconsistencia entre entrenamiento e inferencia (Train-Inference Mismatch). En iteraciones internas posteriores se añadió una corrección para este Mismatch, lo que dio lugar a algoritmos mejorados como GSPO-VR, que aumentan aún más la estabilidad.
\end{warningbox}

\subsection{Diseño de algoritmos a partir de los primeros principios}

Li Yingru (李英儒), desde su perspectiva como investigador teórico del aprendizaje por refuerzo, subrayó dos ideas centrales:

\textbf{Primero, el regreso a los primeros principios.} Muchos de los problemas actuales de colapso en el entrenamiento provienen de descuidar conceptos básicos. Por ejemplo, muchos marcos, en la etapa de Recompute, usan la probabilidad recalculada como si fuera la probabilidad de Rollout, lo cual no coincide con la intuición: la probabilidad que en verdad produjo el muestreo debería provenir del motor de inferencia (como vLLM), no del Recompute del motor de entrenamiento.

\textbf{Segundo, la teoría debe estar al servicio de la práctica.} Li Yingru considera que hacer teoría no es un ejercicio de complejidad por sí misma, sino que sirve para: (1) ofrecer explicaciones simples; (2) proporcionar orientación metodológica para diseñar algoritmos Scalable más simples y efectivos.

\begin{knowledgebox}{El origen histórico del Off-Policy Issue}
El problema del Train-Inference Mismatch es, en esencia, una nueva manifestación en la era de los modelos grandes del clásico Off-Policy Issue. Este problema se ha estudiado durante décadas en el RL tradicional. Li Yingru ha investigado problemas similares en el RL de videojuegos desde 2019 (Divergence of Entity), y el trabajo que hacía entonces es muy parecido al método actual de KL Penalty. Desde TRPO (2015) hasta los diversos métodos de corrección de IS actuales, el núcleo siempre ha sido \textbf{la garantía de una mejora monótona del desempeño}.
\end{knowledgebox}

\subsection{¿Qué características debe tener un «buen algoritmo de RL»?}

Los cuatro panelistas llegaron a un consenso muy uniforme:

\begin{importantbox}{Los criterios de un buen algoritmo de RL (Recipe)}
\begin{enumerate}
    \item \textbf{Simple}: cada Trick debe ser fácil de implementar, de modo que un principiante pueda entender su función con solo verlo
    \item \textbf{Que coincida con la intuición}: el diseño debe tener una motivation clara, no ser una caja negra del tipo «funciona con solo agregarlo»
    \item \textbf{Efectivo}: capaz de entrenar de manera estable en distintas escalas de modelo y tareas, y de producir beneficios reales de desempeño
    \item \textbf{Composicional}: un buen «algoritmo» es, en realidad, un conjunto de Recipes: un objetivo de optimización básico (como Reinforce + IS Correction) más una serie de Tricks validados (PPO Clip, Batch Norm, Language Consistency Penalty, entre otros)
\end{enumerate}
\end{importantbox}

Zheng Chujie señaló, además, que en lugar de hablar de un «buen algoritmo de RL», sería más apropiado hablar de una «buena receta de RL (Recipe)». En la práctica, una vez que se encuentra un conjunto de recetas estable, la gente simplemente empieza a usarla, sin darle un nombre específico a esa combinación de Tricks.

\subsection{Resumen de la sección}

El diseño de algoritmos de RL está pasando de la búsqueda de un único algoritmo disruptivo a la búsqueda de un conjunto de Recipes estable, sencillo y componible. PRIME representa una nueva forma de modelar el Value Model, Reinforce++ refleja la sabiduría del regreso a lo clásico, y GSPO muestra el poder de la derivación a partir de los primeros principios. Los cuatro panelistas coinciden en que la innovación futura de los algoritmos de RL provendrá cada vez más de la optimización conjunta entre Infra y algoritmos, y no de la derivación teórica pura.
